\documentclass[10pt,a4paper]{amsart}
\usepackage[margin=19mm]{geometry}
\usepackage{amsmath,amssymb,mathtools}
\usepackage[scr=rsfs,cal=euler]{mathalfa}
\usepackage{xfrac}
\usepackage[unicode,hidelinks,bookmarks=false]{hyperref}
\newcommand{\ie}{i.e.\ }
\newcommand{\cf}{cf.\ }
\newcommand{\resp}{respectively\ }
\newcommand{\Subsection}{\subsection}
\providecommand{\nobreakdash}{\nobreak\hbox{-}}
\setlength{\emergencystretch}{2em}
\multlinegap0pt
\usepackage[shortlabels]{enumitem}
\usepackage{amssymb}
\usepackage{mathtools}
\usepackage{bm}
\newtheorem{assumption}{Assumption}
\newtheorem{proposition}{Proposition}
\usepackage{cleveref}
\crefname{assumption}{Assumption}{Assumptions}
\crefname{proposition}{Proposition}{Propositions}
\newcommand{\Bern}{\operatorname{Bern}}
\newcommand{\Pcal}{\mathcal P}
\newcommand{\Scal}{\mathcal S}
\newcommand{\X}{\mathcal X}
\newcommand{\dd}{\,\mathrm d}
\newcommand{\dist}{\operatorname{dist}}
\newcommand{\ip}[2]{\left\langle #1,#2\right\rangle}
\title{Tight Weighted Second-Order Asymptotics for the Wyner--Ahlswede--K\"orner Problem Under Regular Posterior Geometry}
\author{Daming Cao}
\date{Source: arXiv:2608.21991v1 (CC BY 4.0)}
\begin{document}
\maketitle

\noindent\emph{Source and licence.} Tight Weighted Second-Order Asymptotics for the Wyner--Ahlswede--K\"orner Problem Under Regular Posterior Geometry. Version arXiv:2608.21991v1 (CC BY 4.0), distributed by arXiv under the Creative Commons Attribution 4.0 International licence, http://creativecommons.org/licenses/by/4.0/, as recorded in the arXiv licence metadata for this exact version and shown on the abstract page https://arxiv.org/abs/2608.21991v1. Full licence notice: \texttt{licenses/arxiv-2608.21991-CC-BY-4.0.txt}.\par\smallskip

\noindent\textit{Editorial scope.} Counted unit: the article's proposition prop:bsc-regularity (source0.tex lines 2647--2650). The article's complete original proof of it is delivered with it (source0.tex lines 2652--2732, one \texttt{proof} environment of 81 lines). No mathematics is written, repaired or re-derived: the statement and the proof are the article's own lines, in the article's own notation, and no retained line is substituted or re-flowed.  Retained uncounted as the source-local context the proof needs: lines 160--163; lines 191--194; lines 279--292; lines 385--397; lines 404--408; lines 2576--2579; lines 2639--2645.  Nothing was omitted from any retained span, so the package carries no omission record.  No retained line contains a citation, so the wrapper carries no bibliography and no bibliography entry.  No retained line contains a figure environment, an image-inclusion command or drawing code, so no image is delivered and the package declares no asset dependency.  Editorial formatting only, in addition: an AMS \texttt{amsart} preamble that restates the article's own declarations and macros used by the retained lines, namely the environments \texttt{assumption}, \texttt{proposition} on separate counters, together with the standard packages those lines need.  The retained environments print in this document as Assumption 1, Proposition 1 in order of appearance, whereas the article prints the same items with the numbering of its own full document.  The complete native source of the article as distributed in the arXiv e-print for this exact version is bundled unchanged as \texttt{sources/208291/source0.tex}.
\begin{equation}
\phi_c(Q):=\min_{P_{U|X}}\phi_c(Q,P_{U|X}).
\label{eq:phi}
\end{equation}

\begin{assumption}[Essential uniqueness of the optimizer]
\label{ass:unique}
There is an open neighborhood $\mathcal O_0$ of $P$ such that, for every $Q\in\mathcal O_0$, the minimum in \eqref{eq:phi} is attained by a finite-alphabet channel $P_{U|X}^{\star,Q}$.  Remove symbols of zero $Q_U^{\star}$-probability and merge any active symbols that induce the same posterior $Q_{X|U=u}^{\star}$.  The resulting reduced optimizing channel is unique up to a permutation of its output symbols.
\end{assumption}

\begin{proposition}[Posterior primal--dual representation]
\label{prop:posterior-dual}
For every full-support $Q$,
\begin{align}
\phi_c(Q)
&=\inf_{\mu:\,\int q\,\mu(\dd q)=Q_X}
\int F_Q(q)\,\mu(\dd q),
\label{eq:posterior-primal}\\
&=\max_{\kappa\in\mathbb R^{\X}}
\left\{\ip{Q_X}{\kappa}:\ip{q}{\kappa}\le F_Q(q),\ \forall q\in\Pcal(\X)\right\}.
\label{eq:posterior-dual}
\end{align}
An optimal measure can be chosen with at most $|\X|$ atoms.
\end{proposition}

\begin{assumption}[Regular posterior geometry]
\label{ass:regular}
Let $\mathcal O\subset\mathcal O_0$ be an open neighborhood of $P$.  In addition to the essential uniqueness in \cref{ass:unique}, assume that:
\begin{enumerate}[label=\textup{(\alph*)}]
\item  After a consistent relabeling, the optimizer has a fixed number $m$ of active symbols, the maps $Q\mapsto\alpha_u(Q)$ and $Q\mapsto q_u(Q)$ are continuously differentiable, and every $q_u(Q)$ lies in the relative interior of $\Pcal(\X)$.
\item The support function $\phi_c$ is twice continuously differentiable on $\mathcal O$, and the contact set is exactly
\begin{equation}
\Scal_Q=\{q_1(Q),\ldots,q_m(Q)\},\qquad Q\in\mathcal O.
\label{eq:finite-contact}
\end{equation}
\item The residual variance is nondegenerate: $\sigma^2(Q)>0$ throughout a possibly smaller neighborhood of $P$.
\end{enumerate}
\end{assumption}

\begin{equation}
\inf_{\substack{Q\in\mathcal C,\ q\in\Pcal(\X):\\
\dist_1(q,\Scal_Q)\ge\eta}}d_Q(q)>0.
\label{eq:strict-exposure}
\end{equation}

\begin{equation}
c:=\frac{\log((1-\beta)/\beta)}{(1-2p)\log((1-a)/a)}.
\label{eq:bsc-slope-main}
\end{equation}

\begin{align}
F_c'(\theta)
&=\log\frac{\theta}{1-\theta}
+ct\log\frac{1-p\star\theta}{p\star\theta}\\
&=2t\operatorname{artanh}(t(1-2\theta))\,[c-c(\theta)].
\label{eq:bsc-Fprime-app}
\end{align}

\begin{proposition}[Local regularity of the BSC posterior optimizer]
\label{prop:bsc-regularity}
For every $p,\beta\in(0,1/2)$ and the slope $c=c(\beta)$ in \eqref{eq:bsc-slope-main}, the symmetric posterior problem has exactly two contacts, $\beta$ and $1-\beta$, and both contacts have positive second derivative.  Moreover, there is a neighborhood of the symmetric joint distribution in which the two contacts and their positive barycentric weights are unique and smooth, the contact set contains no additional points, and the strict-exposure condition in \eqref{eq:strict-exposure} holds uniformly on compact subsets.  Consequently, \cref{ass:regular}(a)--(b) is satisfied locally.
\end{proposition}

\begin{proof}
We verify the claim in three steps.

\emph{1) The symmetric source has exactly two nondegenerate contacts.}
At $c=c(\beta)$, strict monotonicity of $c(\theta)$ and \eqref{eq:bsc-Fprime-app} give
\begin{equation}
F_c'(\theta)<0\quad(0<\theta<\beta),
\qquad
F_c'(\theta)>0\quad(\beta<\theta<1/2).
\end{equation}
Hence $\beta$ is the unique minimizer of $F_c$ on $[0,1/2]$.  Since $F_c(\theta)=F_c(1-\theta)$, the only global minimizers on $[0,1]$ are $\beta$ and $1-\beta$.  Differentiating the stationarity identity $F_{c(\beta)}'(\beta)=0$ with respect to $\beta$ gives
\begin{equation}
F_c''(\beta)
=-(1-2p)\log\frac{1-a}{a}\,c'(\beta)>0.
\label{eq:bsc-positive-hessian}
\end{equation}
Symmetry gives $F_c''(1-\beta)=F_c''(\beta)>0$.

\emph{2) Continue the contacts under a perturbation of the source distribution.}
For a nearby binary joint distribution $Q$, write
\begin{equation}
F_Q^{\mathrm{bin}}(\theta)
:=D(\Bern(\theta)\|Q_X)
+cH\bigl(\Bern(\theta)V_Q\bigr).
\end{equation}
Two posterior values $\theta_0<\theta_1$ are contacts of a common supporting line with slope $\lambda$ precisely when
\begin{align}
\partial_\theta F_Q^{\mathrm{bin}}(\theta_0)&=\lambda,
\label{eq:bsc-contact-system-1}\\
\partial_\theta F_Q^{\mathrm{bin}}(\theta_1)&=\lambda,
\label{eq:bsc-contact-system-2}\\
F_Q^{\mathrm{bin}}(\theta_1)-F_Q^{\mathrm{bin}}(\theta_0)
&=\lambda(\theta_1-\theta_0).
\label{eq:bsc-contact-system-3}
\end{align}
At the symmetric source, $(\theta_0,\theta_1,\lambda)=(\beta,1-\beta,0)$.  The Jacobian of the left sides of \eqref{eq:bsc-contact-system-1}--\eqref{eq:bsc-contact-system-3} with respect to $(\theta_0,\theta_1,\lambda)$ equals, at this point,
\begin{equation}
\begin{bmatrix}
F_c''(\beta)&0&-1\\
0&F_c''(1-\beta)&-1\\
0&0&-(1-2\beta)
\end{bmatrix}.
\label{eq:bsc-contact-jacobian}
\end{equation}
Its determinant is
\begin{equation}
-(1-2\beta)F_c''(\beta)F_c''(1-\beta)\ne0.
\end{equation}
The implicit function theorem therefore supplies unique smooth solutions
\begin{equation}
Q\longmapsto\theta_0(Q),\qquad
Q\longmapsto\theta_1(Q),\qquad
Q\longmapsto\lambda(Q)
\end{equation}
for all $Q$ in a sufficiently small neighborhood of the symmetric source.

The barycentric weights are then forced by
\begin{align}
\alpha_0(Q)\theta_0(Q)+\alpha_1(Q)\theta_1(Q)&=Q_X(1),\\
\alpha_0(Q)+\alpha_1(Q)&=1.
\end{align}
Explicitly,
\begin{align}
\alpha_0(Q)
&=\frac{\theta_1(Q)-Q_X(1)}{\theta_1(Q)-\theta_0(Q)},\\
\alpha_1(Q)
&=\frac{Q_X(1)-\theta_0(Q)}{\theta_1(Q)-\theta_0(Q)}.
\label{eq:bsc-barycentric-weights}
\end{align}
At the symmetric source, both weights equal $1/2$; hence they remain positive and smooth after the neighborhood is reduced if necessary.

\emph{3) Preserve global exposure and identify the optimizer.}
Choose disjoint closed intervals around $\beta$ and $1-\beta$ on which the second derivative is positive.  By \eqref{eq:bsc-positive-hessian} and continuity, the continued contacts remain strict local minima of the supporting gap throughout a sufficiently small source neighborhood.  On the compact complement of those intervals, the symmetric supporting gap is strictly positive because $\beta$ and $1-\beta$ are the only contacts.  Uniform continuity preserves a positive gap on that complement for all nearby $Q$.  Hence the continued supporting line has exactly the two contacts in \eqref{eq:bsc-contact-system-1}--\eqref{eq:bsc-contact-system-3}, and the strict-exposure constant is positive for every fixed distance from them.

Every optimal posterior distribution is supported on these two contacts by \cref{prop:posterior-dual}; the barycenter equation then forces the unique weights in \eqref{eq:bsc-barycentric-weights}.  In this neighborhood,
\begin{equation}
\phi_c(Q)=\alpha_0(Q)F_Q^{\mathrm{bin}}(\theta_0(Q))
+\alpha_1(Q)F_Q^{\mathrm{bin}}(\theta_1(Q)).
\end{equation}
The function $F_Q^{\mathrm{bin}}(\theta)$ is smooth on a full-support neighborhood, and the implicit-function solutions and weights are smooth.  Hence $\phi_c$ is twice continuously differentiable after the neighborhood is reduced if necessary.  The associated supporting line and its contact set also vary smoothly.  This verifies \cref{ass:regular}(a)--(b).
\end{proof}

\end{document}
